\documentclass[11pt,a4paper]{article}

\usepackage[utf8]{inputenc}
\usepackage[T1]{fontenc}
\usepackage[margin=2.3cm]{geometry}
\usepackage{amsmath}
\usepackage{booktabs}
\usepackage{tabularx}
\usepackage{xcolor}
\usepackage[obeyspaces,spaces]{url}
\usepackage{fancyvrb}
\usepackage{enumitem}
\usepackage{array}

\urlstyle{tt}
\renewcommand{\UrlBreaks}{\do\.\do\_\do\/\do\,\do\-\do\=}
% Identifiers (events, files, options): breakable, underscores allowed.
\newcommand{\id}[1]{\path{#1}}

\newenvironment{note}{\par\smallskip\noindent\begin{tabularx}{\linewidth}{|X|}\hline}{\\\hline\end{tabularx}\par\smallskip}

\fvset{fontsize=\small,frame=single,framesep=2mm}
\setlist{itemsep=2pt,topsep=4pt}
\emergencystretch=3em
\renewcommand{\arraystretch}{1.15}

\title{Machine Setup and First Measurements on iccluster113\\
\large Extended report --- CS-471 project ``Tiered Memory Management Beyond Hotness''}
\author{Mayeul Besnard \and Sergio Boffi}
\date{November 10, 2026}

\begin{document}
\maketitle

\begin{abstract}
This report describes the server provided for our reproduction of Liu et al.\
(OSDI'25), what we verified on it, and the first measurements. It extends
\id{Machine_Setup_Report.md} with background explanations and with the work of
November 10: the PEBS sampling test, the changes to the authors' microbenchmark,
and the first run of experiment E1. All results are \textbf{provisional}: they
were taken with the CPU frequency governor, turbo boost and automatic NUMA
balancing left at their original values. The commands we ran are listed, with
their purpose and outcome, in \id{Commands_Log.md}.
\end{abstract}

\tableofcontents
\newpage

% =====================================================================
\section{Background}

This section explains the concepts used in the rest of the report.

\subsection{Tiered memory and its emulation}

A tiered memory system has a small, fast tier (local DRAM) and a larger, slower
tier, such as CXL-attached memory. The operating system decides on which tier
each page lives. Systems such as TPP, NBT, Nomad and Colloid move frequently
accessed (\emph{hot}) pages to the fast tier.

We have no CXL hardware. Like the authors, we emulate the slow tier with the
memory of the \emph{other} socket of a two-socket server. A program runs on the
CPUs of node~0, and its memory is placed either on node~0 (fast tier) or on
node~1 (slow tier). Every access to node~1 crosses the link between the two
sockets, which adds latency and limits bandwidth. The tool \id{numactl} chooses
the CPUs (\id{--cpunodebind}, \id{--physcpubind}) and the memory node
(\id{--membind}) of a whole process. The library call
\id{numa_alloc_onnode} places a single buffer.

\subsection{Memory-level parallelism and AOL}

An out-of-order core does not wait for each cache miss before issuing the next
one. If several misses are independent, they are in flight at the same time and
their latencies overlap. The average number of misses in flight is the
\emph{memory-level parallelism} (MLP).

\begin{itemize}
  \item In a \textbf{pointer chase} each load needs the address returned by the
        previous one. Misses are serial, MLP is about~1, and every access pays
        the whole memory latency.
  \item In a \textbf{sequential read} the addresses are known in advance. Many
        misses overlap and the hardware prefetchers fetch lines early. MLP is
        high, and a slower memory costs much less than its latency suggests.
\end{itemize}

The paper's metric is the \textbf{Amortized Offcore Latency},
$\text{AOL} = \text{Latency}/\text{MLP}$: the latency each miss
\emph{effectively} costs once overlap is taken into account. The paper's main
claim is that the slowdown caused by slow memory depends on AOL, not on how
often data is accessed. A frequently accessed array with high MLP can tolerate
the slow tier better than a rarely accessed linked list.

\subsection{Measuring AOL with hardware counters}

Intel CPUs count off-core requests, i.e.\ requests that leave the core towards
L3 and memory. Three counters are needed:

\begin{center}
\begin{tabularx}{\linewidth}{>{\raggedright\arraybackslash}p{0.43\linewidth}X}
\toprule
Event & Meaning \\
\midrule
\id{offcore_requests.demand_data_rd} ($N$) & number of demand data reads \\
\id{offcore_requests_outstanding.demand_data_rd} ($O$) & every cycle, adds the number of reads in flight \\
\id{offcore_requests_outstanding.cycles_with_demand_data_rd} ($C$) & cycles with at least one read in flight \\
\bottomrule
\end{tabularx}
\end{center}

$O$ is the sum, over all requests, of their latencies (Little's law). Therefore
\[
\text{Latency} = \frac{O}{N}, \qquad
\text{MLP} = \frac{O}{C}, \qquad
\text{AOL} = \frac{\text{Latency}}{\text{MLP}} = \frac{C}{N}.
\]
AOL does not depend on $O$. This turned out to be important, because $O$ is
unreliable on our machine (Section~\ref{sec:anomaly}).

\subsection{The slowdown model}

The paper models the slowdown on the slow tier as $S = P \times K(\text{AOL})$:
\begin{itemize}
  \item $P$ is the fraction of cycles in which the core is stalled waiting for
        memory. A program that rarely waits for memory cannot slow down much.
  \item $K(\text{AOL}) = 1/(a + b/\text{AOL})$ grows with AOL: misses that are
        not hidden by parallelism make the program more sensitive. The
        authors' scripts use $a = 0.87$ and $b = 24.67$, fitted on Skylake. We
        must fit them again on our machine (experiment E2).
\end{itemize}
On Skylake $P$ is measured with \id{cycle_activity.stalls_l3_miss}. This event
does not exist on our CPU (Section~\ref{sec:counters}).

\subsection{Counters, \texttt{perf} and multiplexing}

\id{perf stat} counts events for one command. A CPU has a few fixed counters
and, on our CPU, 4 programmable counters per hardware thread when
hyper-threading is on. If more events are requested than counters are free,
the kernel \emph{multiplexes}: it rotates the events on the counters and scales
the counts. \id{perf} then prints a percentage such as \texttt{(61.27\%)}, the
fraction of time the event was actually counted. Scaled values are estimates,
so we avoid multiplexing altogether.

\subsection{PEBS}

\emph{Precise Event-Based Sampling} makes the CPU write a record every $n$
occurrences of an event, e.g.\ every 3001 loads that miss L3. For memory events
the record contains the address of the instruction and the address of the
\emph{data}. SOAR uses these samples to find out which allocated objects cause
the misses. The record format depends on the CPU generation; from memory, a
timestamp field was added only with format~3 (Skylake).

\subsection{Linux mechanisms that move memory}

\begin{itemize}
  \item \textbf{NUMA balancing} (\id{kernel.numa_balancing} $=1$) periodically
        samples the pages of a process and migrates them towards the node that
        accesses them. It can undo a placement chosen on purpose.
  \item \textbf{Memory tiering mode} ($=2$, used by NBT) promotes hot pages from
        a slower tier to a faster one, and \textbf{demotion} moves cold pages
        down. Both work only between \emph{different memory tiers}, which the
        kernel builds from each node's ``abstract distance''.
  \item \textbf{CPU frequency}: with a load-dependent governor and turbo, the
        clock changes during a run, so the same latency in nanoseconds
        corresponds to a different number of cycles.
\end{itemize}

% =====================================================================
\section{Summary}

\begin{center}
\small
\begin{tabularx}{\linewidth}{>{\raggedright\arraybackslash}p{0.2\linewidth}X>{\raggedright\arraybackslash}p{0.17\linewidth}}
\toprule
Item & Result & Status \\
\midrule
CPU & Intel Xeon E5-2680 v3 (Haswell-EP), 2 sockets $\times$ 12 cores, HT on (48 logical CPUs) & older than the paper's Skylake \\
Memory & 2 NUMA nodes, about 126\,GB each, almost all free & OK \\
Kernel & \id{6.8.0-147-generic} (Ubuntu) & OK \\
\id{perf} & v6.8.12, usable with \id{sudo} & OK \\
AOL counters & all available & OK \\
$P$ counter & \id{stalls_l3_miss} missing; \id{stalls_l2_pending} as proxy & methodology change \\
Counter budget & 4 programmable counters $\Rightarrow$ two runs per measurement & solved \\
Multiplexing & none in any result file (checked on Nov.~9) & OK \\
Frequency & \id{intel_cpufreq}, \id{schedutil}, turbo on & to be fixed (TA) \\
NUMA balancing & on; only 209 pages migrated during E1 & to be turned off (TA) \\
Memory tiers & a single tier with both nodes & NBT/TPP cannot migrate \\
PEBS & format 2, no timestamp; \id{perf record -d} works & SOAR profiling feasible \\
Uncore frequency & driver not available & natural remote latency \\
Slow tier & pointer chasing $\approx$102\,ns local, $\approx$154\,ns remote ($1.5\times$) & less slow than the paper ($2.1\times$) \\
Counter anomaly & $O$ unreliable; AOL $= C/N$ consistent & AOL derived without $O$ \\
GAPBS & \id{bc}, \id{bfs}, \id{pr}, \id{cc}: +29\% to +37\% on the remote node & provisional \\
Microbenchmark & three changes for E1, verified & done \\
E1 & first run done; whole-program time still missing & provisional \\
SoarAlto artifact & full SOAR pipeline, but Skylake-specific; baselines are kernel patches & porting needed \\
\bottomrule
\end{tabularx}
\end{center}

% =====================================================================
\section{CPU and memory topology}

\begin{Verbatim}
$ lscpu | grep -E "Model name|Vendor|Socket|Core\(s\)|Thread\(s\)|NUMA|^CPU\(s\)|L3"
CPU(s):                 48
Model name:             Intel(R) Xeon(R) CPU E5-2680 v3 @ 2.50GHz
Thread(s) per core:     2
Core(s) per socket:     12
Socket(s):              2
L3 cache:               60 MiB (2 instances)
NUMA node0 CPU(s):      0,2,4,...,46   (even)
NUMA node1 CPU(s):      1,3,5,...,47   (odd)

$ numactl -H
node 0 size: 128686 MB   free: 127073 MB
node 1 size: 129012 MB   free: 128267 MB
node distances:   0: 10 21   1: 21 10
\end{Verbatim}

\begin{itemize}
  \item \textbf{12 cores per socket} are more than the 8 threads used in the paper.
  \item \textbf{30\,MiB of L3 per socket} is far smaller than the workloads
        (several GB), so they really exercise memory.
  \item \textbf{About 126\,GB per node}: any workload fits entirely on one node,
        which is what static placement experiments need.
  \item \textbf{CPU numbering is interleaved}: even CPUs are on node~0, odd CPUs
        on node~1. A range such as \id{0-11} would spread threads over both
        sockets. This is exactly the bug we found in the microbenchmark
        (Section~\ref{sec:changes}).
  \item \textbf{Hyper-thread siblings}: CPU $i$ and CPU $i+24$ share a physical
        core. CPUs \id{0,2,4,6,8,10,12,14} are 8 distinct physical cores of
        node~0; we use them for multi-threaded runs.
  \item The NUMA distance 21 is an estimate written by the firmware, not a
        measurement. The real latencies are in Section~\ref{sec:micro}.
\end{itemize}

% =====================================================================
\section{Performance counters}
\label{sec:counters}

\subsection{Access}

\id{/proc/sys/kernel/perf_event_paranoid} is \texttt{4}, a Debian/Ubuntu value
that forbids counters to normal users. We run \id{sudo perf} for each
measurement, instead of lowering this setting, so the machine configuration
stays unchanged.

\subsection{Events and derived metrics}

All AOL events exist. For $P$ we use \id{cycle_activity.stalls_l2_pending},
the cycles stalled while an L2 miss is pending.

\begin{center}
\begin{tabular}{ll}
\toprule
Quantity & Formula \\
\midrule
Latency (cycles) & $O / N$ \\
MLP & $O / C$ \\
AOL (cycles) & $C / N$ \\
$P$ & \id{stalls_l2_pending} / \id{cycles} \\
\bottomrule
\end{tabular}
\end{center}

\textbf{Differences from the paper}, due to Haswell:
\begin{itemize}
  \item \id{stalls_l2_pending} also counts stalls on L2 misses that hit in L3.
        Our $P$ therefore includes stalls the slow tier does not affect, and is
        probably overestimated.
  \item There are no offcore events restricted to L3 misses. Latency and AOL mix
        L3 hits (about 40 cycles) with DRAM accesses, so AOL can be ``diluted''
        for cache-friendly workloads.
  \item A possible correction, to be evaluated in E2, uses Intel's Top-Down
        method for Haswell. From memory, and still to be checked against
        Intel's documentation:
        \[
        P_\text{DRAM} \approx \Bigl(1 - \frac{h}{h + 7m}\Bigr)
        \times \frac{\text{stalls\_l2\_pending}}{\text{cycles}},
        \]
        with $h$ = \id{mem_load_uops_retired.l3_hit} and
        $m$ = \id{mem_load_uops_retired.l3_miss}. It would need a third run
        per measurement.
\end{itemize}

\subsection{Multiplexing and the two-run procedure}

A first test with six events showed multiplexing on \id{cycles} (61\%) and
\id{stalls_l2_pending} (39\%). The cause is the NMI watchdog
(\id{nmi_watchdog} $=1$): it uses the fixed cycle counter, so our
\id{cycles} must take a programmable one, and with hyper-threading only four
are available. Five programmable events do not fit.

Instead of disabling the watchdog or hyper-threading, which would change the
machine, each measurement is split into two runs of the same command, both
verified to run without multiplexing:

\begin{Verbatim}
# Run A - AOL
sudo perf stat -e cycles,instructions,offcore_requests.demand_data_rd,\
offcore_requests_outstanding.demand_data_rd,\
offcore_requests_outstanding.cycles_with_demand_data_rd -- <cmd>
# Run B - P
sudo perf stat -e cycles,instructions,cycle_activity.stalls_l2_pending -- <cmd>
\end{Verbatim}

The price is that $P$ and AOL come from different executions; with stable runs
(standard deviation below 1\%) this is acceptable. On November~9 we checked all
result files for multiplexing percentages: there are none.

% =====================================================================
\section{CPU frequency}

\begin{Verbatim}
scaling_driver: intel_cpufreq      scaling_governor: schedutil      no_turbo: 0
cpuinfo_min/max_freq: 1200000 / 3300000 kHz
available governors: conservative ondemand userspace powersave performance schedutil
\end{Verbatim}

\id{intel_cpufreq} is \id{intel_pstate} in passive mode, so the generic
governors can be used. With \id{schedutil} and turbo, the frequency moves
between 1.2 and 3.3\,GHz depending on load and temperature. For final
measurements we plan, after TA approval, the \id{performance} governor and turbo
off (fixed 2.5\,GHz), restoring the original values afterwards.

% =====================================================================
\section{Kernel tiering state}

\begin{Verbatim}
$ uname -r; cat /proc/sys/kernel/numa_balancing; cat /sys/kernel/mm/numa/demotion_enabled
6.8.0-147-generic
1
false
$ cat /sys/devices/virtual/memory_tiering/memory_tier*/nodelist
0-1                       (a single tier: memory_tier4)
\end{Verbatim}

\begin{itemize}
  \item \textbf{NUMA balancing is on.} It could move our pages during a run.
        During E1 the system-wide counter \id{numa_pages_migrated} grew by
        only 209 pages, against about 500\,000 pages in our buffers, so it did
        not matter there. It should still be off for final runs.
  \item \textbf{Both nodes are in the same memory tier.} Promotion and demotion
        only move pages between tiers, so NBT and TPP would never migrate
        anything. Taking node~1's CPUs offline, as the authors do, is probably
        not enough, since the tier depends on the abstract distance set by the
        memory driver (to be verified). The artifact contains a module,
        \id{run/colloid/tpp/tierinit/tierinit.c}, that moves node~1 to a lower
        tier, but it relies on helpers defined in the Colloid kernel patch.
\end{itemize}

\begin{center}
\begin{tabular}{ll}
\toprule
Experiment & Works on the stock kernel? \\
\midrule
E1 (static placements) & yes (\id{numactl}, \id{numa_alloc_onnode}) \\
E2 (counters, calibration of $K$) & yes \\
SOAR profiling & yes, with \id{perf record -d} (Section~\ref{sec:pebs}) \\
SOAR placement (\id{LD_PRELOAD} interception) & yes by code review; not tested \\
E3 baselines NBT, TPP & no: two memory tiers needed \\
E3 baselines Nomad, Colloid & no: custom kernels \\
\bottomrule
\end{tabular}
\end{center}

% =====================================================================
\section{The authors' artifact (SoarAlto)}

\subsection{Structure}

\begin{center}
\small
\begin{tabularx}{\linewidth}{>{\raggedright\arraybackslash}p{0.25\linewidth}X}
\toprule
Path & Content \\
\midrule
\id{src/microbenchmark/} & pointer-chasing and sequential microbenchmark (E1) \\
\id{src/soar/} & SOAR: allocation tracing (\id{prof/ldlib.c}), sampling (\id{run/prof.sh}), object scoring (\id{run/proc_obj_e.py} $\rightarrow$ \id{obj_stat.csv}), placement by interception (\id{interc/ldlib.c}, policy in \id{check_trace}), patches (\id{nbt.patch}, \id{gapbs.patch}) \\
\id{src/alto/} & kernel patches and build scripts for TPP, NBT, Nomad, Colloid, with and without ALTO \\
\id{run/} & \id{setup.sh}, \id{config.sh}, and \id{bc-urand/} (the only workload with a run directory) \\
\bottomrule
\end{tabularx}
\end{center}

\subsection{Findings that matter on our machine}

\begin{enumerate}
  \item \textbf{SOAR does not need the authors' kernel.} The SOAR README asks for
        Linux 5.18 with \id{nbt.patch}. By inspection, the only part of the
        patch SOAR uses is a \id{pebs:pebs} tracepoint (instruction address,
        data address, timestamp); the rest concerns NBT migrations. Our PEBS
        format has no timestamp, so the tracepoint would be useless anyway, and
        \id{perf record -d} replaces it (Section~\ref{sec:pebs}).
  \item \textbf{Slow-tier emulation.} The authors lower the uncore frequency of
        the remote socket to 500\,MHz (about 190\,ns latency) through
        \id{/sys/devices/system/cpu/intel_uncore_frequency}. This driver is not
        available on our machine. The alternative is writing MSR \id{0x620},
        which needs TA approval; otherwise we keep the natural remote latency.
  \item \textbf{Local DRAM capacity is limited with the GRUB \id{memmap=}
        parameter}, a persistent boot change that conflicts with the TA rule
        that the stock kernel stays the default. Reserving 1\,GB huge pages on
        node~0 is a possible alternative (untested).
  \item \textbf{Baselines} are built by scripts that switch kernel version,
        patch, install and update the bootloader. They must be adapted to the
        one-time GRUB boot procedure. ``Colloid'' in the artifact is
        Colloid-TPP, with a Skylake-specific patch.
  \item \textbf{Scripts must be read before running.} \id{src/soar/run/prof.sh}
        calls \id{check_cxl_conf}, which disables the NMI watchdog, ASLR, KSM,
        NUMA balancing, THP, hyper-threading, turbo and swap and takes node~1's
        CPUs offline. It must not be run as is.
\end{enumerate}

\subsection{SOAR on Haswell}

\begin{center}
\small
\begin{tabularx}{\linewidth}{>{\raggedright\arraybackslash}p{0.27\linewidth}X}
\toprule
Phase & Work required \\
\midrule
Allocation tracing (\id{prof/ldlib.c}) & change the hard-coded output path; may need a fix for recent \id{g++} (\id{throw(std::bad_alloc)}) \\
Sampling (\id{run/prof.sh}) & replace the tracepoint with \id{perf record -d} (done in principle); align the \id{perf} clock with the \id{rdtsc} timestamps of \id{ldlib.c}; fix hard-coded paths \\
Object scoring (\id{proc_obj_e.py}) & Haswell event names; recalibrated $a$, $b$ and AOL thresholds (40, 45, 60, 80) \\
Placement (\id{interc/ldlib.c}) & replace the call-site addresses in \id{check_trace} with ours \\
Workloads & GAPBS built with \id{-no-pie} (done); Intel MLC not available \\
\bottomrule
\end{tabularx}
\end{center}

% =====================================================================
\section{System settings and tools}

\begin{Verbatim}
THP: always [madvise] never      ASLR: 2      KSM: 0      SMT: on
GRUB_DEFAULT=0   GRUB_TIMEOUT_STYLE=hidden   GRUB_TIMEOUT=0   empty command lines
1 GB huge pages on node 0: 0
df: 162 GB free; /, /boot and /home on the same partition
\end{Verbatim}

\begin{itemize}
  \item \textbf{THP} is \id{madvise}: plain \id{malloc} does not get huge pages.
        Our pointer chase therefore also suffers TLB misses with 4\,KB pages.
  \item \textbf{ASLR} is on. With \id{-no-pie}, the call-site addresses SOAR
        relies on should stay fixed, so ASLR may not need to be disabled.
  \item \textbf{GRUB}: \id{GRUB_DEFAULT=0} boots the first menu entry, and the
        menu is hidden with no timeout. A newly installed custom kernel with a
        higher version could become the default, and could not be avoided by
        hand at boot. This must be fixed before any custom kernel.
  \item \textbf{Tools}: \id{gcc}, \id{make} and \id{libnuma-dev} were missing.
        We installed \id{build-essential} and \id{libnuma-dev} with \id{apt}.
\end{itemize}

% =====================================================================
\section{First measurements with the microbenchmark}
\label{sec:micro}

\textbf{Conditions:} unmodified microbenchmark, one thread on node~0, the whole
process's memory on node~0 or node~1 (\id{numactl --membind}), 1\,GiB buffer.

\subsection{Quick timings}

\begin{center}
\begin{tabular}{lrrr}
\toprule
Pattern (\id{-i 3}) & Node 0 (s) & Node 1 (s) & Slowdown \\
\midrule
Pointer chasing & 15.2 & 22.0 & +45\% \\
Sequential read & 8.9 & 11.6 & +30\% \\
\bottomrule
\end{tabular}
\end{center}

The low-MLP pattern slows down more, as the paper predicts. These times include
the construction and shuffling of the list.

\subsection{Counters (\texttt{-i 9}, 3 runs per configuration)}

\begin{center}
\small
\begin{tabular}{lrrrrrr}
\toprule
Configuration & Time (s) & Latency & MLP & AOL & $P$ & Slowdown \\
\midrule
Pointer chasing, node 0 & 25.49 & 603 & 3.7 & 162 & 0.76 & --- \\
Pointer chasing, node 1 & 37.01 & 988 & 4.2 & 236 & 0.79 & +45\% \\
Sequential, node 0 & 23.46 & 183 & 11.7 & 16 & 0.66 & --- \\
Sequential, node 1 & 31.42 & 252 & 12.3 & 21 & 0.74 & +34\% \\
\bottomrule
\end{tabular}
\end{center}

Latency and AOL are in cycles. $P$ is similar for the two patterns, while AOL
differs by a factor of~10, which is the kind of difference the model needs.
Latency and MLP in this table are not reliable (Section~\ref{sec:anomaly}).

\subsection{Pointer chasing without the setup phase}

The setup phase is the same for any number of iterations. Subtracting a run with
\id{-i 1} from the run with \id{-i 9} leaves 8 traversals, about 134 million
steps:

\begin{center}
\begin{tabular}{lrr}
\toprule
Quantity (\id{-i 9} minus \id{-i 1}) & Node 0 & Node 1 \\
\midrule
Time per step & \textbf{$\approx$102\,ns} & \textbf{$\approx$154\,ns} \\
Core cycles per step & $\approx$253 & $\approx$383 \\
$C$ per step & $\approx$254 & $\approx$385 \\
Requests per step & 1.35 & 1.38 \\
\bottomrule
\end{tabular}
\end{center}

The remote node is about \textbf{1.5$\times$} slower, against about 2.1$\times$
(90 vs 190\,ns) in the paper. The time per step includes TLB misses, so it is
not a pure memory latency. Our slow tier is therefore milder than the paper's,
and this must be stated in the methodology.

\subsection{Counter anomaly}
\label{sec:anomaly}

A pointer chase is serial, so its MLP should be close to~1. We measure 4.2--5.2,
also after removing the setup phase, and a latency of 791 cycles while one step
takes 253 cycles: a request cannot be outstanding for longer than the step that
waits for it. Hence \textbf{$O$ (\id{offcore_requests_outstanding.demand_data_rd})
is not reliable on this machine}, at least for this pattern. The cause is
unknown (a Haswell erratum or an effect of hyper-threading are possible).

$C$ matches the real cycles per step, so \textbf{AOL $= C/N$ remains usable}.
The authors' \id{proc_obj_e.py} computes it the same way. Latency and MLP
derived from $O$ will not be reported for pointer chasing.

% =====================================================================
\section{First GAPBS results}

\textbf{Conditions:} GAPBS built with \id{-no-pie}; graph \id{urand24}
(16.8\,M vertices, 2.2\,GB) read once into the page cache from node~0; 8 OpenMP
threads on cores \id{0,2,...,14}; memory on node~0 or node~1; 3 runs with each
event set. Arguments: \id{bc -i4 -n1}, \id{pr -n1}, \id{bfs -n8}, \id{cc -n8}.
\id{sssp} needs a weighted graph and \id{tc} is very slow, so they were not
run. Summary: \id{gapbs_summary.log} in the repository.

\begin{center}
\small
\begin{tabular}{lrrrrll}
\toprule
Workload & Local (s) & Remote (s) & Slowdown & Whole process & $P$ (loc./rem.) & AOL (loc./rem.) \\
\midrule
\id{bc}  & 17.32 & 23.75 & \textbf{+37.1\%} & +35.2\% & 0.78 / 0.84 & 46 / 63 \\
\id{bfs} & 0.24  & 0.33  & \textbf{+35.4\%} & +25.2\% & 0.56 / 0.64 & 35 / 48 \\
\id{pr}  & 4.29  & 5.74  & \textbf{+33.8\%} & +30.1\% & 0.87 / 0.90 & 28 / 38 \\
\id{cc}  & 0.59  & 0.76  & \textbf{+29.1\%} & +23.6\% & 0.65 / 0.72 & 62 / 82 \\
\bottomrule
\end{tabular}
\end{center}

The slowdown uses GAPBS's \emph{Average Time}, which covers the algorithm only.

\begin{itemize}
  \item The slowdowns are close to each other (29--37\%), partly because our
        slow tier is only 1.5$\times$ slower. A narrow range leaves little for a
        predictor to explain.
  \item The correlation between local $P$ and slowdown is weak (about 0.2), but
        four points say nothing about the hypotheses yet.
  \item The counters cover the whole process, including graph loading. For
        \id{bfs} and \id{cc} the algorithm is a small part of the run, so their
        $P$ and AOL mostly describe loading. More trials (\id{-n}) are needed.
\end{itemize}

% =====================================================================
\section{PEBS sampling test (November 10)}
\label{sec:pebs}

\textbf{Question:} can we collect what SOAR needs (L3-miss samples with data
address and time) on the stock kernel?

\begin{Verbatim}
sudo perf record -d -e mem_load_uops_retired.l3_miss:P -c 3001 \
     -o ~/hotness/results/pebs_test.data -- \
     numactl --cpunodebind=0 --membind=0 ./bench -R 0.0 -i 1 -A 1024
\end{Verbatim}

\id{:P} requests precise sampling (PEBS), \id{-c 3001} takes one sample every
3001 events as the authors do, and \id{-d} records the data address.

\textbf{Result:} 17\,814 samples. Where they fall:

\begin{center}
\begin{tabular}{lr}
\toprule
Function & Samples \\
\midrule
\id{shuffle} (list setup) & 33.4\% \\
\id{pc_thread} (traversal; \id{ptr_chase} is inlined here) & 31.4\% \\
\id{init_ptr_buf} (list setup) & 31.1\% \\
kernel & $\approx$4\% \\
\bottomrule
\end{tabular}
\end{center}

Samples inside \id{pc_thread} look like this (time in seconds, data address,
instruction address, function):

\begin{Verbatim}
1894399.201677:     749a8bdbfc80     5c6c278f00e6 pc_thread
1894399.201993:     749a88f4f640     5c6c278f00e6 pc_thread
1894399.202305:     749a89de5080     5c6c278f00e6 pc_thread
\end{Verbatim}

Three independent checks show that the samples are correct:
\begin{enumerate}
  \item The data addresses differ at every sample and are multiples of 64: they
        are list nodes, one cache line each. The instruction address is always
        the same load of the loop.
  \item The addresses lie within about 57\,MB, consistent with the list being
        shuffled in blocks of about 64\,MB.
  \item Samples are about 0.316\,ms apart, i.e.\ 3001 misses, or
        \textbf{$\approx$105\,ns per miss}: the same as the 102\,ns per step
        measured from execution times.
\end{enumerate}

\textbf{Conclusions:}
\begin{itemize}
  \item \textbf{SOAR profiling is feasible on this machine}, on the stock
        kernel, without the authors' patch or Linux 5.18.
  \item The timestamps come from the kernel's \id{perf} clock (time since boot),
        not from the PEBS record. The remaining work is to align them with the
        \id{rdtsc} timestamps written by \id{ldlib.c}, e.g.\ by switching
        \id{ldlib.c} to the same clock or by selecting a clock with the
        \id{-k} option of \id{perf record} (to be tested).
  \item With one iteration, two thirds of the misses come from the setup phase.
        This explains why whole-program counters were dominated by setup.
\end{itemize}

% =====================================================================
\section{Microbenchmark changes for E1 (November 10)}
\label{sec:changes}

E1 runs two threads at the same time: one chases pointers in buffer~A (a
1\,GiB linked list, low MLP), the other reads buffer~B sequentially (a 1\,GiB
array, high MLP). The array is the \emph{hot} data: with \id{-i 9} it is read
$9 \times 26 = 234$ times, about 3.9 billion cache-line accesses, against about
151 million for the list. The unmodified program could not run this experiment.
We made three changes on branch \id{haswell-port} of our clone of SoarAlto.

\begin{enumerate}
  \item \textbf{Pinning on one socket.} Thread $i$ was pinned to CPU $i$, so the
        two threads ran on CPUs 0 and 1, which are on different sockets here.
        Now thread $i$ is pinned to CPU $2i$. Check: the two worker threads run
        on CPUs 0 and 2, both on node~0.
  \item \textbf{One memory node per buffer.} Both buffers came from plain
        \id{malloc}; the existing function \id{init_buf} (which calls
        \id{numa_alloc_onnode}) was never used, and option \id{-r} had no
        effect. Both threads now use \id{init_buf} and free with
        \id{numa_free}; \id{-r} selects the node of buffer~A, and a new option
        \id{-b} selects the node of buffer~B. Buffer~B is allocated 4096 bytes
        larger, because \id{bandwidth()} writes 8 bytes past its end, which
        could crash once the buffer ends exactly at a page boundary. Check with
        \id{numastat}: about 1025\,MB on node~0 and 1024\,MB on node~1 with
        \id{-r 0 -b 1}.
  \item \textbf{Per-thread timing.} The program printed nothing. Each thread now
        measures only its work loop with \id{clock_gettime(CLOCK_MONOTONIC)}
        and prints \id{pc_time_s} or \id{bw_time_s}. Check: with \id{-i 1},
        \id{pc_time_s} is 1.72\,s, i.e.\ 102\,ns $\times$ 16.8\,M steps, as
        expected.
\end{enumerate}

% =====================================================================
\section{E1: first results (November 10, provisional)}

Four placements, 5 repetitions each, \id{-i 9}, both buffers 1\,GiB. The names
refer to the hotness-based view, in which the array is the hot data.

\begin{center}
\small
\begin{tabular}{llrrrr}
\toprule
Configuration & List / array & \multicolumn{2}{c}{Pointer chasing} & \multicolumn{2}{c}{Sequential} \\
 & node & time (s) & slowdown & time (s) & slowdown \\
\midrule
\id{all_dram} & 0 / 0 & 15.74 & --- & 22.62 & --- \\
\id{all_slow} & 1 / 1 & 24.16 & +53.5\% & 30.94 & +36.8\% \\
\id{hot_on_dram} & 1 / 0 & 23.28 & +47.8\% & 22.69 & +0.3\% \\
\id{cold_on_dram} & 0 / 1 & 16.06 & +2.0\% & 30.85 & +36.4\% \\
\bottomrule
\end{tabular}
\end{center}

\textbf{The placement stayed static:} \id{numa_pages_migrated} grew by 209 pages
in the whole system during the 20 runs.

\textbf{Reading the results:}
\begin{itemize}
  \item \textbf{In relative terms, \id{cold_on_dram} wins.} Placing the cold but
        latency-bound list in DRAM costs the two threads about 38\% in total,
        against 48\% for the hotness-based choice. This is the direction of
        H1.
  \item \textbf{The sequential read suffers more than expected (+36\%).} It
        reads about 11\,GB/s from node~0 and 8\,GB/s from node~1
        (234\,GiB in 22.6 and 30.9\,s). On this machine the remote array seems
        limited by the bandwidth of the link between the sockets, not only by
        latency, so its MLP protects it less than in the paper.
  \item \textbf{In seconds lost, the two placements are close:} about 7.6\,s for
        \id{hot_on_dram} and 8.5\,s for \id{cold_on_dram}, because the array
        thread does far more work.
  \item \textbf{The decisive metric is missing:} the time of the whole program.
        The program ends with the slower thread, and the list thread also spends
        about 10\,s in its setup, which accesses memory too. A run with
        whole-program time (3 repetitions) is the next step.
\end{itemize}

% =====================================================================
\section{Original state of the machine}

To be restored after any experiment that changes it.

\begin{center}
\small
\begin{tabular}{ll}
\toprule
Setting & Original value \\
\midrule
\id{scaling_governor} & \id{schedutil} \\
\id{intel_pstate/no_turbo} & \id{0} \\
\id{scaling_min_freq} / \id{scaling_max_freq} & 1200000 / 3300000 \\
\id{kernel.numa_balancing} & 1 \\
\id{numa/demotion_enabled} & false \\
\id{kernel.nmi_watchdog} & 1 \\
\id{kernel.perf_event_paranoid} & 4 \\
\id{transparent_hugepage/enabled} & \id{madvise} \\
\id{kernel.randomize_va_space} & 2 \\
\id{ksm/run} & 0 \\
\id{smt/control} & \id{on} \\
1\,GB huge pages on node 0 & 0 \\
\id{GRUB_DEFAULT} / \id{GRUB_TIMEOUT} / \id{GRUB_TIMEOUT_STYLE} & 0 / 0 / \id{hidden} \\
Kernel command line additions & none \\
Kernel & \id{6.8.0-147-generic} \\
\bottomrule
\end{tabular}
\end{center}

\textbf{Changes made so far:} packages \id{build-essential} and \id{libnuma-dev};
files under \id{~/hotness/}; a repository-local git identity in our clone of
SoarAlto. Possibly also the \id{msr} module and \id{msr-tools}, suggested by
another chat; whether they were run is unconfirmed (see \id{Commands_Log.md}).
No system setting has been changed.

% =====================================================================
\section{Open questions for the TA}

\begin{enumerate}
  \item During measurements, may we set the governor to \id{performance},
        disable turbo and set \id{numa_balancing} to 0, restoring the values
        right afterwards?
  \item Both NUMA nodes are in one memory tier, so NBT and TPP cannot migrate.
        Is there a recommended way to make node~1 a lower tier?
  \item The artifact limits local DRAM with the GRUB \id{memmap=} parameter. Is
        that allowed (e.g.\ for a one-time boot only), or should we use another
        method, such as 1\,GB huge pages?
  \item The uncore frequency driver is not available. May we set the remote
        uncore frequency through MSR \id{0x620}, or should we keep the natural
        remote latency (1.5$\times$)?
  \item Which of the artifact's settings (HT, swap, ASLR, KSM, THP, NMI
        watchdog, node~1's CPUs offline) may we apply during our runs?
  \item We installed \id{build-essential} and \id{libnuma-dev}. Is that fine,
        and may we also install \id{vmtouch} and \id{libgfortran5}?
  \item Could we get CloudLab access (c220g5, the authors' Skylake platform) for
        the baselines?
  \item Who else uses this machine, and is a remote console available if a
        kernel does not boot?
\end{enumerate}

% =====================================================================
\section{Next steps}

\begin{enumerate}
  \item E1 with whole-program time; copy \id{e1.log} to the repository and plot.
  \item Repeat E1 and the counter measurements under stable settings, after TA
        approval.
  \item E2: more varied workloads (more GAPBS trials, a pointer chase with $N$
        independent chains to vary MLP in a controlled way), then fit $a$ and
        $b$.
  \item SOAR: align the \id{perf} and \id{ldlib.c} clocks, port the scoring
        script, then try one workload.
  \item Record the commit hashes of SoarAlto and GAPBS.
  \item Before any custom kernel: persistent journal, \id{/boot} space check,
        one-time GRUB boot.
\end{enumerate}

\end{document}
